\documentclass[11pt]{article}
\input{../../shared/project_style.tex}
\rhead{Basics 13 Textbook Draft}

\begin{document}
\DocTitle{Basics 13: Lorentz Force and Charged-Particle Motion}{II - Electromagnetism foundations}{Gyromotion, helical trajectories, electric acceleration, drift motion, and numerical integration}

\begin{overviewbox}
\textbf{Textbook draft, version 1.0.} The Lorentz equation is the microscopic starting point for orbit following, guiding-center theory, and kinetic descriptions. This chapter solves the simplest field configurations exactly and uses them to develop gyrofrequency, Larmor radius, pitch angle, and the electric-cross-magnetic drift. This chapter is designed for a reader with no prior plasma-physics training. It distinguishes exact identities from physical approximations and connects the central equations to later simulation modules.
\end{overviewbox}

\ProjectTOC

\section{Learning objectives}
After completing this note, the reader should be able to:
\begin{itemize}[leftmargin=1.6em]
\item Solve charged-particle motion in uniform electric and magnetic fields.
\item Derive gyrofrequency and Larmor radius.
\item Separate parallel and perpendicular motion.
\item Derive the electric-cross-magnetic drift.
\item Explain why drift velocity is independent of particle mass and charge sign.
\item Assess numerical orbit integration using invariants.
\end{itemize}

\section{Prerequisites and reading strategy}
\begin{itemize}[leftmargin=1.6em]
\item Basics 2, 4, and 11.
\end{itemize}
On a first reading, emphasize physical meaning and dimensional checks. On a second reading, reproduce the derivations. Attempt the computational laboratory only after the limiting cases are understood.

\section{Equation of motion and energy balance}

For a nonrelativistic particle,
\begin{equation}
m\frac{d\mathbf v}{dt}=q\left(\mathbf E+\mathbf v\times\mathbf B\right),
\qquad
\frac{d\mathbf x}{dt}=\mathbf v.
\end{equation}
Dotting with $\mathbf v$ gives
\begin{equation}
\frac{d}{dt}\left(\frac{mv^2}{2}\right)=q\mathbf v\cdot\mathbf E.
\end{equation}
Electric fields can change kinetic energy. Magnetic fields change direction but not speed. This separation is a powerful diagnostic for orbit codes.


\section{Uniform magnetic field: gyromotion}

Let $\mathbf B=B_0\hat{\mathbf z}$ and $\mathbf E=0$. Then
\begin{align}
\dot v_x&=\Omega_s v_y,\\
\dot v_y&=-\Omega_s v_x,\\
\dot v_z&=0,
\end{align}
where the signed cyclotron frequency is
\begin{equation}
\Omega_s=\frac{qB_0}{m}.
\end{equation}
Differentiating once more gives $\ddot v_x+\Omega_s^2v_x=0$. The perpendicular velocity rotates at $|\Omega_s|$, while $v_\parallel=v_z$ is constant. The Larmor radius is
\begin{equation}
\rho_L=\frac{v_\perp}{|\Omega_s|}=\frac{mv_\perp}{|q|B_0}.
\end{equation}
The sign of $q$ determines the direction of rotation.


\section{Helical trajectories and pitch angle}

Integrating velocity produces a circle in the plane perpendicular to $\mathbf B$ and uniform translation along $\mathbf B$. The result is a helix. Define pitch angle $\alpha$ by
\begin{equation}
v_\parallel=v\cos\alpha,\qquad v_\perp=v\sin\alpha.
\end{equation}
The distance advanced along the field during one gyroperiod $T_c=2\pi/|\Omega_s|$ is
\begin{equation}
L_{\rm pitch}=v_\parallel T_c.
\end{equation}
Do not confuse this geometric helix pitch with pitch variables used for magnetic-mirror or guiding-center coordinates.


\section{Uniform electric field}

With $\mathbf B=0$ and constant $\mathbf E$,
\begin{equation}
\mathbf v(t)=\mathbf v_0+\frac{q\mathbf E}{m}t,\qquad
\mathbf x(t)=\mathbf x_0+\mathbf v_0t+\frac{q\mathbf E}{2m}t^2.
\end{equation}
For $\mathbf E=E_\parallel\hat{\mathbf b}$ in the presence of a uniform magnetic field, the parallel motion accelerates while perpendicular gyromotion persists. A sustained parallel electric field therefore changes both energy and pitch angle.


\section{Crossed fields and the $\mathbf E\times\mathbf B$ drift}

Let $\mathbf E\perp\mathbf B$ be constant. Seek a constant velocity $\mathbf v_E$ for which the transverse Lorentz force vanishes:
\begin{equation}
\mathbf E+\mathbf v_E\times\mathbf B=0.
\end{equation}
Cross with $\mathbf B$ and use $\mathbf E\cdot\mathbf B=0$:
\begin{equation}
\boxed{\mathbf v_E=\frac{\mathbf E\times\mathbf B}{B^2}}.
\end{equation}
Writing $\mathbf v=\mathbf v_E+\mathbf w$ reduces the equation for $\mathbf w$ to ordinary gyromotion. The guiding center drifts at $\mathbf v_E$, independent of $m$, $|q|$, and charge sign. Hence ions and electrons share this bulk drift and it does not itself create a current in a uniform plasma.


\section{Ordering toward guiding-center theory}

If field scales $L$ and times $\tau$ satisfy
\begin{equation}
\epsilon_\rho=\frac{\rho_L}{L}\ll1,\qquad
\epsilon_\Omega=\frac{1}{|\Omega_s|\tau}\ll1,
\end{equation}
the rapid gyrophase can be separated from slow guiding-center motion. Nonuniform fields then produce grad-$B$, curvature, polarization, and other drifts. Those are developed in later kinetic notes.


\section{Numerical orbit integration}

A naive explicit Euler step can increase particle energy in a static magnetic field. Suitable checks include
\begin{equation}
\Delta K,\qquad \Delta v_\parallel,\qquad
\Delta \rho_L,\qquad \text{phase error per gyroperiod}.
\end{equation}
The Boris algorithm is widely used because its magnetic rotation preserves speed to high accuracy. Adaptive Runge-Kutta methods may be accurate locally but do not automatically preserve long-time geometric structure.


\section{Computational laboratory}
Implement analytic and numerical trajectories for uniform $\mathbf B$, parallel $\mathbf E$, and crossed $\mathbf E$ and $\mathbf B$. Compare Euler, Runge-Kutta, and Boris updates over hundreds of gyroperiods. Measure energy, phase, guiding-center, and Larmor-radius errors.

\section{Summary}
\begin{itemize}[leftmargin=1.6em]
\item The Lorentz equation separates electric work from magnetic redirection.
\item Uniform magnetic fields produce gyromotion characterized by gyrofrequency and Larmor radius.
\item Parallel motion combines with gyromotion to form a helix.
\item Crossed uniform fields produce a species-independent $\mathbf E\times\mathbf B$ drift.
\item Orbit solvers should be judged by invariant and phase preservation, not only local truncation error.
\end{itemize}

\SymbolsAcronymsSection
\begin{symtable}
$\Omega_s$ & signed cyclotron frequency & $qB/m$. \\
$\rho_L$ & Larmor radius & $v_\perp/|\Omega_s|$. \\
$v_\parallel,v_\perp$ & velocity components & Parallel and perpendicular to the field. \\
$\alpha$ & pitch angle & Angle between velocity and magnetic field. \\
$\mathbf v_E$ & electric-cross-magnetic drift & $\mathbf E\times\mathbf B/B^2$. \\
$T_c$ & gyroperiod & $2\pi/|\Omega_s|$. \\
\end{symtable}

\section{Exercises}
\begin{enumerate}[leftmargin=1.8em]
\item Derive the gyromotion solution including the charge-sign dependence.
\item Show that the magnetic moment $mv_\perp^2/(2B_0)$ is constant in a uniform field.
\item Derive the $\mathbf E\times\mathbf B$ drift using a shifted velocity.
\item Calculate electron and deuteron gyrofrequencies at a specified magnetic field.
\item Explain why shared $\mathbf E\times\mathbf B$ motion produces no current in a uniform quasineutral plasma.
\end{enumerate}

\section{References}
\begin{thebibliography}{99}
\bibitem{ref1} G. B. Arfken, H. J. Weber, and F. E. Harris, \emph{Mathematical Methods for Physicists}, 7th ed., Academic Press (2013).
\bibitem{ref2} G. Strang, \emph{Linear Algebra and Its Applications}, 4th ed., Brooks/Cole (2006).
\bibitem{ref3} W. A. Strauss, \emph{Partial Differential Equations: An Introduction}, 2nd ed., Wiley (2007).
\bibitem{ref4} D. J. Griffiths, \emph{Introduction to Electrodynamics}, 4th ed., Cambridge University Press (2017).
\bibitem{ref5} J. D. Jackson, \emph{Classical Electrodynamics}, 3rd ed., Wiley (1998).
\bibitem{ref6} F. F. Chen, \emph{Introduction to Plasma Physics and Controlled Fusion}, 3rd ed., Springer (2016).
\bibitem{ref7} R. J. Goldston and P. H. Rutherford, \emph{Introduction to Plasma Physics}, CRC Press (1995).
\bibitem{ref8} P. M. Bellan, \emph{Fundamentals of Plasma Physics}, Cambridge University Press (2006).
\end{thebibliography}

\end{document}
